% sec:fd-stencils — Finite Difference Stencil Derivation
\section{Finite Difference Stencil Derivation}
\label{sec:fd-stencils}

A grid stores numbers at discrete points; the differential equations of
numerical relativity demand derivatives at those same points.  The gap between
the two is bridged by a single tool: Taylor expansion.  Expanding a smooth
function about a grid point and algebraically eliminating the function value in
favour of its derivatives produces a \emph{finite difference stencil} --- a
recipe for approximating a derivative from a fixed pattern of neighbouring grid
values.  The order of the stencil determines how rapidly the approximation
improves as the grid is refined, and the derivation itself is mechanical enough
to be automated for arbitrary order.

\paragraph{Setup and notation.}

Consider a uniform grid with spacing $\Delta x$ and grid points
$x_i = x_0 + i\,\Delta x$.  Let $f$ be a smooth function sampled on this
grid, with $f_i = f(x_i)$.  The goal is to approximate the $p$-th derivative
$f^{(p)}(x_i) = \d^p\!f/\d x^p\big|_{x_i}$ using the values
$f_{i+k}$ at a finite set of offsets $k \in \{k_1, k_2, \ldots, k_N\}$
from the point of interest.  A centred stencil uses offsets symmetric about
$k = 0$; a one-sided stencil places all offsets on one side of the target
point.
\physnote{The uniform-spacing assumption is essential: non-uniform grids
require modified stencils or coordinate mappings.  The codebase uses uniform
grids within each refinement level, so the standard stencils apply directly.}

\paragraph{Taylor expansion.}

The key identity is the Taylor expansion of the shifted value
$f_{i+k} = f(x_i + k\,\Delta x)$ about $x_i$:
%
\begin{multline}
  f_{i+k} = \sum_{n=0}^{\infty} \frac{(k\,\Delta x)^n}{n!}\,f^{(n)}(x_i) \\
  = f_i + k\,\Delta x\,f'_i + \frac{k^2\,\Delta x^2}{2}\,f''_i
    + \frac{k^3\,\Delta x^3}{6}\,f'''_i + \cdots
  \label{eq:taylor-grid}
\end{multline}
%
Each shifted value is an infinite power series in $k\,\Delta x$ with
coefficients built from the successive derivatives of $f$ at $x_i$.  The derivative we want ---
say $f'_i$ --- appears as the coefficient of $k\,\Delta x$.  If we can choose
weights $w_k$ such that $\sum_k w_k\,f_{i+k}$ cancels the unwanted
derivatives and isolates $f'_i$, we have our stencil.

\paragraph{The Vandermonde system.}

To make this precise, suppose we seek weights $\{w_k\}$ for offsets
$k \in \{k_1, \ldots, k_N\}$ such that the linear combination
$\sum_j w_{k_j}\,f_{i+k_j}$ approximates
$\Delta x^p\,f^{(p)}_i$.  Substituting
\cref{eq:taylor-grid} for each $f_{i+k_j}$ and collecting powers of
$\Delta x$ gives the conditions
%
\begin{equation}
  \sum_{j=1}^{N} w_{k_j}\,\frac{k_j^n}{n!}
  = \begin{cases}
      1 & \text{if } n = p \,, \\
      0 & \text{if } n \neq p \,,
    \end{cases}
  \qquad n = 0, 1, \ldots, N-1 \,.
  \label{eq:stencil-conditions}
\end{equation}
%
Each row of this system kills one unwanted derivative: with $N$ points we
have $N$ degrees of freedom, enough to zero out $N - 1$ unwanted terms and
set the coefficient of the desired derivative to one.  The coefficient
matrix has $(n, j)$ entry $k_j^n/n!$; up to row scaling by $n!$, this is a
Vandermonde matrix, and it is nonsingular whenever the offsets $k_j$ are
distinct.
\histnote{The Vandermonde system for finite difference coefficients dates to
Fornberg's 1988 algorithm \cite{Fornberg1988}, which generates stencils of
arbitrary order and derivative with $\order{N^2}$ operations.}

The system always has a unique solution.  What varies with the choice of
offsets is the \emph{order of accuracy}: the power of $\Delta x$ at which the
first uncancelled derivative appears.  With $N$ stencil points, we can satisfy
$N$ conditions in \cref{eq:stencil-conditions} and cancel derivatives through
order $n = N - 1$.  The leading error term is then proportional to
$\Delta x^{N-p}$, since the first uncancelled term has $n = N$ and contributes
at order $\Delta x^N$ before dividing by $\Delta x^p$ to recover the
derivative.

\paragraph{Centred stencils.}

The most efficient stencils for interior points are centred: they use offsets
$k \in \{-m, \ldots, -1, 0, 1, \ldots, m\}$, giving $N = 2m + 1$ stencil
points.  Symmetry about $k = 0$ forces the weights of odd-powered terms to
cancel in pairs when approximating an even-order derivative, and vice versa.
For first and second derivatives --- the two cases needed throughout this book
--- a centred stencil with $2m + 1$ points achieves order $2m$.  For even
derivatives such as $f''$, the symmetry cancellation provides one extra order
of accuracy beyond what a one-sided stencil with the same number of points
would achieve.

For the first derivative ($p = 1$) with $m = 1$ (three points:
$k \in \{-1, 0, 1\}$), the conditions \cref{eq:stencil-conditions} reduce to
%
\begin{equation}
  w_{-1} + w_0 + w_1 = 0 \,, \qquad
  -w_{-1} + w_1 = 1 \,, \qquad
  \tfrac{1}{2}\,w_{-1} + \tfrac{1}{2}\,w_1 = 0 \,.
  \label{eq:second-order-d1-system}
\end{equation}
%
The first equation cancels $f_i$ (zeroth derivative); the second isolates
$\Delta x\,f'_i$; the third cancels $\Delta x^2\,f''_i/2$.  The solution is
$w_{-1} = -1/2$, $w_0 = 0$, $w_1 = 1/2$, giving the familiar second-order
centred difference
%
\begin{equation}
  f'_i \approx \frac{f_{i+1} - f_{i-1}}{2\,\Delta x} \,,
  \label{eq:d1-second-order}
\end{equation}
%
with a leading error term proportional to $\Delta x^2\,f'''_i / 6$.  The error
is second order in $\Delta x$ despite using only three points --- a
consequence of the antisymmetry of the weights, which automatically cancels the
even-powered terms.

\paragraph{Higher-order stencils.}

Extending to $m = 2$ (five points: $k \in \{-2, -1, 0, 1, 2\}$) and solving
the $5 \times 5$ Vandermonde system yields the fourth-order centred first
derivative:
%
\begin{equation}
  f'_i \approx \frac{-f_{i+2} + 8\,f_{i+1} - 8\,f_{i-1} + f_{i-2}}
  {12\,\Delta x} \,.
  \label{eq:d1-fourth-order}
\end{equation}
%
The leading error is now $\order{\Delta x^4}$: doubling the resolution reduces
the error by a factor of sixteen.  The same five-point width produces a
fourth-order second derivative:
%
\begin{equation}
  f''_i \approx \frac{-f_{i+2} + 16\,f_{i+1} - 30\,f_i + 16\,f_{i-1}
    - f_{i-2}}{12\,\Delta x^2} \,.
  \label{eq:d2-fourth-order}
\end{equation}
%
The large negative central weight ($-30$) in the second derivative stencil
reflects the operator's sensitivity to curvature: small changes in $f_i$
relative to its neighbours produce large second-derivative estimates, which
is why second derivatives amplify grid-scale noise more aggressively than
first derivatives.  Both stencils appear directly in the codebase; comparing
the numerator weights $(-1, 8, -8, 1)$ and $(-1, 16, -30, 16, -1)$ with
the Rust implementation confirms exact correspondence.
\coderef{stencils}{rs}

\paragraph{Mixed partial derivatives.}

The BSSN evolution equations require mixed second derivatives
$\partial^2 f/\partial x\,\partial y$.  The natural two-dimensional extension
of the five-point stencil is a $5 \times 5$ grid of $25$ points, and
accessing $25$ scattered memory locations per grid point per derivative can
dominate the cost of a stencil evaluation.  A more efficient approach exploits
the fact that mixed partials commute for smooth functions: apply the
one-dimensional first derivative stencil in $x$ at each of the five
$y$-offsets, then apply the $y$-stencil to the five intermediate results.
Both passes access the same $25$ grid points, but the separable structure
improves cache behaviour --- data is read in structured row patterns rather
than a scattered two-dimensional block --- and simplifies the implementation
to two reuses of the one-dimensional stencil routine.
\implnote{The codebase function \texttt{d2xy} implements the two-pass
strategy: first computing $\partial f/\partial x$ at five $y$-neighbours,
then applying the $y$-stencil to the five intermediate results.  The combined
denominator is $144\,\Delta x\,\Delta y = (12)^2\,\Delta x\,\Delta y$.}

The separable approach preserves the fourth-order accuracy of the individual
stencils.  To see why, note that each one-dimensional pass introduces an error
of $\order{\Delta x^4}$ (or $\order{\Delta y^4}$).  The total error of the
mixed derivative approximation is therefore
$\order{\Delta x^4 + \Delta y^4}$, the same order as the unmixed stencils.
On a uniform grid with $\Delta x = \Delta y$, the mixed derivative
converges at the same rate as the pure derivatives.

\paragraph{Boundary stencils.}

Centred stencils require data on both sides of the target point, which is
unavailable at grid boundaries.  The standard remedy is a one-sided stencil
that shifts all offsets to the interior.  For a second-order forward
approximation to the first derivative, the offsets are $k \in \{0, 1, 2\}$
and the Vandermonde system gives
%
\begin{equation}
  f'_i \approx \frac{-3\,f_i + 4\,f_{i+1} - f_{i+2}}{2\,\Delta x} \,,
  \label{eq:d1-onesided}
\end{equation}
%
with a leading error of $\order{\Delta x^2}$.  The backward version
(offsets $k \in \{0, -1, -2\}$) has the same coefficients with reversed
signs.  The drop from fourth-order interior accuracy to second-order boundary
accuracy is typical: one-sided stencils need more points to achieve the same
order as centred stencils, and boundary regions are usually a small fraction
of the total grid.
\warnnote{One-sided boundary stencils \emph{can} reduce the global convergence
order --- for hyperbolic problems, a boundary scheme of order $q - 1$
typically preserves interior order $q$, but lower-order closures may degrade
the global rate.  In practice, the boundary treatment is often dominated by
other considerations --- outgoing wave conditions, excision, or AMR buffer
zones --- and the second-order boundary stencil is adequate.}

\paragraph{Stencil width and ghost zones.}

A centred stencil with $2m + 1$ points reaches $m$ grid points on each side
of the target.  The integer $m$ is the stencil's \emph{half-width}.  For the
fourth-order stencils used in the codebase, $m = 2$: the stencil spans five
points, and the function \texttt{d1x} can only be evaluated at interior points
satisfying $2 \leq i < N_x - 2$, where $N_x$ is the number of grid points in
the $x$ direction.  Rather than switching to one-sided stencils at every
boundary, we can extend the grid by $m$ extra points on each side ---
\emph{ghost zones} --- and fill them with data from boundary conditions or
neighbouring patches in an adaptive mesh refinement hierarchy
(\cref{ch:amr}).  The centred stencil then applies uniformly at every
interior point, simplifying the inner loop.

The Kreiss--Oliger dissipation operator (\cref{sec:fd-dissipation}) uses a
wider seven-point stencil with half-width $m = 3$, requiring one additional
layer of ghost zones.  The codebase respects the maximum stencil width when
allocating ghost zones, ensuring that all stencils can be applied at every
interior point without bounds violations.
